\section{Daten bearbeiten mit \texorpdfstring{\gls{SQL}}{SQL}}
Mit SQL kann man Daten nicht nur abfragen, sondern auch erstellen, verändern und löschen. Im folgenden Schauen wir uns dafür nötigen Befehle an.

\subsection{Eine neue Tabelle erstellen: \texorpdfstring{\lstinline[style=sql]|CREATE table|}{CREATE table}}
\label{subsec:create-table}
Eine neue Tabelle kann innerhalb einer Datenbank mit dem Befehl \lstinline[style=sql]|CREATE TABLE| angelegt werden:

\begin{lstlisting}[style=sql]
CREATE TABLE tabellenname(
  kolonne1 eigenschaften1,
  kolonne2 eigenschaften2,
  ...
)
\end{lstlisting}

\begin{myexample}
\label{ex:sql-create}
Die Tabelle \texttt{users} wurde mit folgendem Befehl erstellt:

\begin{lstlisting}[style=sql]
CREATE TABLE users (
  `id` int(10) UNSIGNED NOT NULL AUTO_INCREMENT PRIMARY KEY,
  `username` varchar(191) NOT NULL UNIQUE,
  `email` varchar(191) NOT NULL UNIQUE,
  `password` varchar(191) NOT NULL,
  `name` varchar(191) NOT NULL,
  `bio` varchar(191) DEFAULT NULL,
  `gender` enum('male','female') DEFAULT NULL,
  `birthday` datetime DEFAULT NULL,
  `city` varchar(191) DEFAULT NULL,
  `country` varchar(191) DEFAULT NULL,
  `centimeters` int(11) DEFAULT NULL,
  `avatar` varchar(191) NOT NULL DEFAULT 'avatar.png',
  `role` enum('user','dba','teacher','admin') NOT NULL DEFAULT 'user',
  `is_active` tinyint(1) NOT NULL DEFAULT 0,
  `remember_token` varchar(100) DEFAULT NULL,
  `created_at` timestamp NULL DEFAULT NULL,
  `updated_at` timestamp NULL DEFAULT NULL
)
\end{lstlisting}
\end{myexample}

Nach dem Kolonnennamen folgt zuerst der \textbf{Kolonnentyp}, d.h., die Angabe, welche Art von Daten in der Kolonne gespeichert werden. Einige der häufigsten Datentypen sind:
\begin{itemize}
\item \lstinline[style=sql]|INT(laenge)| bezeichnet eine Ganzzahl, für die \texttt{laenge} bytes zur Verfügung stehen. Beispielsweise speichert die Kolonne \texttt{centimeters} eine Zahl von 0 bis $2^{11}$, d.h. 2048.
\item \lstinline[style=sql]|TINYINT(1)| speichert den Wert 0 oder 1, wobei 0 typischerweise ``falsch'' und 1 ``wahr'' bedeuten.
\item \lstinline[style=sql]|DATETIME| ist ein Datum mit einer Uhrzeit
\item \lstinline[style=sql]|VARCHAR(laenge)| bezeichnet eine Zeichenkette, d.h. ein Text, der mit maximal \texttt{laenge} vielen Bit kodiert wird.
\item \lstinline[style=sql]|ENUM('wert1', 'wert2',...)| bezeichnet eine Auflistung (engl. \textit{enumeration}) möglicher Werte. Beispielsweise können nur ``male'' und ``female'' im Feld \texttt{gender} eingetragen werden.
\end{itemize}

Zudem können folgende Eingeschaften für Kolonnen angegeben werden:

\begin{itemize}
	\item \lstinline[style=sql]|NOT NULL|: Feld darf nicht leer sein. Falls kein \lstinline[style=sql]|DEFAULT|-Wert angegeben wird, muss, der Wert für diese Kolonne beim Erstellen eines neuen Benutzerkontos immer mitgeliefert werden.
	\item \lstinline[style=sql]|DEFAULT wert|: Standard-Wert, falls nichts anderes angegeben wird. Das Attribut \texttt{avatar} nimmt beispielsweise den Wert \texttt{avatar.png} an, falls kein anderes Bild hochgeladen wird.
	\item \lstinline[style=sql]|AUTO_INCREMENT|: Zahl, die bei jedem neuen Eintrag automatisch immer um 1 grösser wird
	\item \lstinline[style=sql]|PRIMARY KEY|: Schlüsselattribut (s. \cref{table:table-cia}).
	\item \lstinline[style=sql]|UNIQUE|: Darf keine doppelten Werte enthalten
\end{itemize}

\subsection{Tabellen löschen: \texorpdfstring{\lstinline[style=sql]|DROP table|}{DROP table}}
Mit dem Befehl \lstinline[style=sql]|DROP TABLE tabellenname| kann eine Tabelle komplett löschen. 

\begin{myachtung}
Dieser Schritt ist irreversibel. Falls Sie eine Tabelle irrtümlich löschen, geben Sie mir Bescheid, damit ich Ihre Datenbank zurücksetzen kann.
\end{myachtung}

\begin{aufgabe}{aufg:a-sql-insert}
Erstellen Sie eine Tabelle \texttt{users2} mit denselben Eigenschaften wie \texttt{users} (s. Beispiel \ref{ex:sql-create})
\begin{myantwort}
\begin{lstlisting}[style=sql]
CREATE TABLE users2 (
  `id` int(10) UNSIGNED NOT NULL AUTO_INCREMENT PRIMARY KEY,
  `username` varchar(191) NOT NULL UNIQUE,
  `email` varchar(191) NOT NULL UNIQUE,
  `password` varchar(191) NOT NULL,
  `name` varchar(191) NOT NULL,
  `bio` varchar(191) DEFAULT NULL,
  `gender` enum('male','female') DEFAULT NULL,
  `birthday` datetime DEFAULT NULL,
  `city` varchar(191) DEFAULT NULL,
  `country` varchar(191) DEFAULT NULL,
  `centimeters` int(11) DEFAULT NULL,
  `avatar` varchar(191) NOT NULL DEFAULT 'avatar.png',
  `role` enum('user','dba','teacher','admin') NOT NULL DEFAULT 'user',
  `is_active` tinyint(1) NOT NULL DEFAULT 0,
  `remember_token` varchar(100) DEFAULT NULL,
  `created_at` timestamp NULL DEFAULT NULL,
  `updated_at` timestamp NULL DEFAULT NULL
)
\end{lstlisting}
\end{myantwort}
\end{aufgabe}

\begin{aufgabe}{}
Löschen Sie danach die Tabelle \texttt{users2} wieder.
\begin{myantwort}
\begin{lstlisting}[style=sql]
DROP TABLE users2
\end{lstlisting}
\end{myantwort}
\end{aufgabe}


\subsection{Daten einfügen: \texorpdfstring{\lstinline[style=sql]|INSERT|}{INSERT}}

Mit \lstinline[style=sql]|INSERT| können neue Zeilen in eine bestehende Tabelle eingefügt werden.

Die Syntax von \lstinline[style=sql]|INSERT| sieht wie folgt aus:

\begin{lstlisting}[style=sql]
INSERT INTO tabelle_name (kolonne1, kolonne2, kolonne3, ...)
VALUES (wert1, wert2, wert3, ...)
\end{lstlisting}

\begin{myexample}[sql-insert]
Folgender Code fügt eine neue Person in die Tabelle \texttt{users} ein:
\begin{lstlisting}[style=sql]
INSERT INTO users (username, email, password, name, bio, gender, birthday, city, country, centimeters, avatar, role, is_active, remember_token, created_at, updated_at) 
VALUES ('guenther37', 'guenther@instahub.test', '12345', 'Günther Müller', 'Günther mag Kartoffelsalat.', 'male', '2006-06-06 00:00:00', 'Leipzig', 'Deutschland', '173', 'avatar.png', 'user', '0', NULL, now(), now())
\end{lstlisting}
Wie Sie sehen, wird zuerst die Tabelle (\texttt{users}) aufgelistet, danach die Kolonnen und schliesslich die Werte.
\end{myexample}

\begin{aufgabe}{aufg:sql-insert}
Für welche Kolonnen müssen für einen neuen Benutzer immer die Werte mitgegeben werden? S. die Erklärungen zu \lstinline[style=sql]|NOT NULL| und \lstinline[style=sql]|DEFAULT| in \cref{subsec:create-table}.
\begin{myantwort}
\begin{verbatim}
username, email, password, name, role, is_active
\end{verbatim}
\end{myantwort}
\end{aufgabe}
\begin{myabgabe}
Erstellen Sie einen neuen Eintrag in der Tabelle \texttt{users} mit folgenden Eigenschaften:
\begin{table}[H]
\centering
\begin{tabular}{c|c}
\textbf{Kolonne} & \textbf{Wert} \\\midrule
\texttt{username} & \texttt{testuser}\\
\texttt{email} & \texttt{test@test.com} \\
\texttt{passwort}& \texttt{test123}\\
\texttt{name}& \texttt{Test-Vorname Test-Nachname}\\
\texttt{role} & \texttt{user}\\
\texttt{is\_active} & \texttt{1}\\
\end{tabular}
\end{table}
Überprüfen Sie danach, ob der Eintrag korrekt erstellt worden ist, indem Sie  \lstinline[style=sql]|SELECT * FROM users| ausführen.
Ihr(e) Tisch-Nachbar(in) soll nun versuchen, sich mit diesen Angaben auf Ihrer InstaHub-Seite einzuloggen.
\begin{myantwort}
\begin{lstlisting}[style=sql]
INSERT INTO users (username, email, password, name, role, is_active) 
VALUES ('testuser', 'test@test.com', 'test123', 'Test-Vorname Test-Nachname', 'user', '1')
\end{lstlisting}
\end{myantwort}
\end{myabgabe}

\subsection{Daten verändern: \texorpdfstring{\lstinline[style=sql]|UPDATE|}{UPDATE}}

Mit \lstinline[style=sql]|UPDATE| können Werte einer Tabelle aktualisiert werden. Die Syntax von \lstinline[style=sql]|UPDATE| sieht wie folgt aus:

\begin{lstlisting}[style=sql]
UPDATE tabelle_name
SET kolonne1 = wert1, kolonne2=wert2, ...
WHERE bedingung(en)
\end{lstlisting}

\begin{aufgabe}{aufgabe:evalPoly}

Ersetzen Sie die Stadt ``Berlin'' überall durch ``Bern''
\begin{myantwort}
\begin{lstlisting}[style=sql]
UPDATE users
SET city="Bern"
WHERE city="Berlin"
\end{lstlisting}
\end{myantwort}
\end{aufgabe}

\begin{myabgabe}[hjdecounter]
 Setzen Sie die Körpergrössen aller männlichen Mitglieder auf 190.
\begin{myantwort}
\begin{lstlisting}[style=sql]
UPDATE users 
SET centimeters=190 
WHERE gender="male"
\end{lstlisting}
\end{myantwort}
\end{myabgabe}

\subsection{Einträge löschen: \texorpdfstring{\lstinline[style=sql]|DELETE FROM|}{DELETE FROM}}
Mit dem Befehl \lstinline[style=sql]|DELETE FROM| können Einträge aus einer Tabelle gelöscht werden.
\begin{lstlisting}[style=sql]
DELETE FROM tabellenname
WHERE bedingungen
\end{lstlisting}

\begin{myexample}
Folgender Code löscht den Eintrag für user \texttt{guenther37}:
\begin{lstlisting}[style=sql]
DELETE FROM users
WHERE username="guenther37"
\end{lstlisting}
\end{myexample}

\begin{aufgabe}{}
Löschen Sie den Eintrag, den Sie in \cref{aufg:a-sql-insert} erstellt haben. Überprüfen Sie, ob der Eintrag verschwunden ist, indem Sie danach \lstinline[style=sql]|SELECT * FROM users| ausführen.
%TODO Referenz stimmt nicht
\begin{myantwort}
\begin{lstlisting}[style=sql]
DELETE FROM users
WHERE username="testuser"
\end{lstlisting}
\end{myantwort}
\end{aufgabe}

\section{Weiterführende Links}
Folgende Links dienen der Vertiefung in das Thema SQL (und der Prüfungsvorbereitung).

Allgemein zu Datenbanken: 
\begin{itemize}
\item \url{https://oinf.ch/kurs/vernetzung-und-systeme/datenbanken/}
\item \url{https://wi-wissen.github.io/instahub-doc-de/#/exercices}
\end{itemize}
Zu \gls{SQL}:
\begin{itemize}
\item \url{https://www.w3schools.com/sql/default.asp} (auf Englisch)
\item \url{https://sql-island.informatik.uni-kl.de}
\item \url{https://sql-tutorial.de/home/lektionen.php?lektion=1}. Hier wird unter anderem folgende Datenbank verwendet:
\begin{figure}[H]
\centering
\begin{tikzpicture}[entity/.style={ellipse, draw=black, minimum width=2cm}, key/.style={rectangle, draw=black, minimum width=2cm, minimum height=1cm}]

% Entity "laender"
\node[key] (laender) at (0,0) {cia};

% Attribute "Name" with underlining to denote a key
\node[entity, left=of laender] (name) {\underline{name}};
\draw (laender) -- (name);

% Attribute "Region"
\node[entity, above left=of laender] (region) {region};
\draw (laender) -- (region);

% Attribute "Fläche"
\node[entity, above=of laender] (flaeche) {flaeche};
\draw (laender) -- (flaeche);

% Attribute "Einwohner"
\node[entity, above right=of laender] (einwohner) {einwohner};
\draw (laender) -- (einwohner);

% Attribute "BIP"
\node[entity, right=of laender] (bip) {bip};
\draw (laender) -- (bip);

\end{tikzpicture}
\end{figure}
\end{itemize}


